\documentclass[tikz,border=8pt]{standalone}
\usepackage{ctex}
\usepackage{amsmath}
\usetikzlibrary{calc, arrows.meta}

\begin{document}
\begin{tikzpicture}[
  scale=1.0, thick,
  vec/.style={-{Stealth[length=7pt,width=4pt]}, thick},
]

%% ====== 左图：平行四边形法则 ======

% 公共起点 O
\coordinate (O) at (0,0);

% F1 方向（右偏上）
\coordinate (F1) at (2.4, 0.7);
% F2 方向（右偏上更多）
\coordinate (F2) at (1.0, 2.0);
% 合力 F = F1 + F2
\coordinate (F)  at (3.4, 2.7);

% 平行四边形虚线
\draw[gray, dashed, thin] (F1) -- (F);
\draw[gray, dashed, thin] (F2) -- (F);

% F1（红色）
\draw[vec, red, thick] (O) -- (F1) node[below right, font=\small] {$\vec{F_1}$};

% F2（蓝色）
\draw[vec, blue, thick] (O) -- (F2) node[left, font=\small] {$\vec{F_2}$};

% 合力 F（黑色加粗）
\draw[vec, black, very thick, line width=1.6pt]
  (O) -- (F) node[above right, font=\small] {$\vec{F}$};

% 起点点
\filldraw[black] (O) circle (2pt) node[below left, font=\small] {$O$};

% 公式
\node[font=\small] at (1.7, -0.8)
  {$\vec{F}=\vec{F_1}+\vec{F_2}$};

%% ====== 右图：船过河 ======
% 右图偏移 (5.5, 0)

\begin{scope}[xshift=5.5cm]

% 河岸（蓝色水面背景）
\fill[blue!10] (-0.4, 0) rectangle (2.8, 2.6);
\draw[blue!50, thin] (-0.4, 0) -- (2.8, 0);
\draw[blue!50, thin] (-0.4, 2.6) -- (2.8, 2.6);
\node[font=\footnotesize, blue!70] at (-0.3, 1.3) {河};

% 公共起点（船位置）
\coordinate (B) at (0.4, 0.5);

% 船速（蓝色，垂直河面）
\coordinate (Vs) at (0.4, 2.0);  % 船速终点，在 B 正上方
\coordinate (Vw) at (1.5, 0.5);  % 水速终点（水平向右）
\coordinate (Vr) at (1.5, 2.0);  % 合速度终点

% 平行四边形辅助线
\draw[gray, dashed, thin] (Vs) -- (Vr);
\draw[gray, dashed, thin] (Vw) -- (Vr);

% 水速（红色）
\draw[vec, red, thick] (B) -- (Vw) node[below, font=\footnotesize] {$\vec{v}_{\text{水}}$};

% 船速（蓝色）
\draw[vec, blue, thick] (B) -- (Vs) node[left, font=\footnotesize] {$\vec{v}_{\text{船}}$};

% 合速（黑粗）
\draw[vec, black, very thick, line width=1.6pt]
  (B) -- (Vr) node[right, font=\footnotesize] {$\vec{v}$};

% 船起点
\filldraw[black] (B) circle (2pt);

% 标题
\node[font=\footnotesize] at (1.1, -0.25) {船过河示意};

\end{scope}

\end{tikzpicture}
\end{document}
